\section{Speed, energy and what the instruments cost}\label{sec:efficiency}

\subsection{A smaller build is not a faster build at the full window}
\begin{table}[h]
\centering\small
\caption{Decode at the full 262{,}144 window (prefill to 94.8\,\%), first battery, MTP $n{=}2$,
\flag{q4\_0} KV, official non-thinking sampling, 192-token generations. Groups match on build, context,
split and checkpoint setting. Spread is (max$-$min)/median; the adjusted column divides each reading by
its tokens-per-pass (\cref{sec:variance}).}\label{tab:fullwindow}
\shrinktofit{\begin{tabular}{@{}llrlrrr@{}}
\toprule
build & split & $n$ & readings (tok/s) & median & spread & spread, adjusted \\
\midrule
\Qfive & 54,46 & 3 & 10.82, 12.70, 12.94 & 12.70 & 16.7\,\% & 2.9\,\% \\
\Qsix & 58,42 & 3 & 10.75, 11.90, 14.16 & 11.90 & 28.6\,\% & 9.7\,\% \\
\Qsix & 56,44 & 3 & 10.22, 11.71, 14.99 & 11.71 & 40.7\,\% & 11.2\,\% \\
\Qfour & 56,44 & 1 & 12.10 & --- & n/a & --- \\
\bottomrule
\end{tabular}}
\end{table}

At the full window the three builds that reach it decode within 7\,\% of each other
(\cref{tab:fullwindow}), while repetitions of one configuration spread by up to 41\,\%. The common
expectation that a smaller quantization buys speed does not hold here: at depth the attention read over
the cache dominates each pass, and it is the same size for every build. The second battery agrees
under cleaner conditions: with MTP at about 100K tokens \Qsix{} and \Qfour{} decode at 16.3 and
17.4\tps; per-pass time at 262K orders the builds by size (184, 191, 194\,ms) with overlapping
ranges. The smaller build earns its place through memory (a gentler cache at the same window,
\cref{sec:kvmap}), not through throughput. That table went through three versions in my notes (32.9\,\%,
46.7\,\%, 40.7\,\% for the same spread) before one was computed with a stated estimator over matched
cells; see \cref{sec:corrections}.

\subsection{Power and energy}
Over a representative 12-hour window of the first battery (26\,\% GPU-busy; 43{,}182 one-second
samples), the two GPUs drew 1.08\,kWh and the CPU package 0.17\,kWh. A modelled system total, adding a
fixed platform allowance, is 1.79\,kWh: 149\,W mean, 80\,W median, 409\,W peak, 334\,W mean under GPU
load. GPU0 peaked at 90\,°C and GPU1 at 76\,°C. The manual splits put more layers on GPU0, which is the
likely cause, but airflow between the slots was not controlled, so this is an association. A split
chosen for context also selects a thermal operating point.

Per-token energy exists only for decode into an empty cache\hist{} (\cref{tab:backends}, and in the
same series 3.7\,J/token for the smallest build with DFlash2 up to 11.8 for \QsixXL{} without
speculation): speculation cuts energy per token by $2$--$2.5\times$ at every build, because the cards
draw similar power for fewer seconds. At depth, the second battery recorded GPU energy per
\emph{request}: a median 37--41\,kJ per AA-LCR answer with DFlash2 against 45--51\,kJ with MTP, and
28--35\,kJ against 66--78\,kJ on GPQA's long reasoning chains. These include prefix-cache effects and
whatever else the host was doing, and 1\,Hz sampling cannot resolve decode bursts, so they are
per-request figures only. A controlled joules-per-token curve against depth was planned and cancelled.

\subsection{What each instrument bought}\label{sec:cost}

\begin{figure}[t]
\centering
\begin{tikzpicture}
\begin{axis}[paper,xbar,height=6.4cm,width=0.70\linewidth,bar width=8pt,
  xmin=0,xmax=98,xlabel={GPU-hours},
  symbolic y coords={aalcr,ruler,he,hs,kldtask,kld},ytick=data,
  yticklabels={{AA-LCR + GPQA, 4 configs (second battery)},{RULER, 12 cells to 131K},{HumanEval+ paired, 2 builds},{HellaSwag $n{=}400$, 4 builds},{KL divergence, task prompts, 4 cells},{KL divergence, prose + code, 10 cells}},
  yticklabel style={font=\scriptsize},
  nodes near coords,nodes near coords style={font=\scriptsize,anchor=west},
  point meta=explicit symbolic,enlarge y limits=0.12]
\addplot[fill=cbgrey!55,draw=black!60] coordinates {
 (68.9,aalcr)[69 h: no separation]
 (12.34,ruler)[12.3 h: saturated / budget effect]
 (1.29,he)[1.3 h: bound of $\pm$3 pts]
 (0.54,hs)[0.5 h: 1-pt spread, 4-bit on top]
 (0.37,kldtask)[0.4 h: all builds ordered]
 (1.95,kld)[2.0 h: all builds ordered, 30--32 of 32 windows]};
\end{axis}
\end{tikzpicture}
\caption{What a few hours of the same machine bought from each instrument. Wall-clock including model
loads, single runs (the second-battery bar sums per-answer generation time: 12 h for AA-LCR, 57 h for
GPQA Diamond). Divergence ranked the ladder in 2.3 GPU-hours; about 83 hours of task benchmarks
across five instrument classes bounded it and ranked nothing. SWE-bench (four days, exploratory period)
predates timing instrumentation and is not shown.}\label{fig:cost}
\end{figure}

\Cref{fig:cost} is the economics of the thesis. The comparison is loose (fixed token budgets against
fixed problem counts, single runs), and it answers different questions: divergence says \emph{how far
the distribution moved}, a benchmark says \emph{whether outcomes changed}. For choosing between
neighbouring quantizations of one model, the first is cheap and decisive, and the second is expensive
and, at public benchmark sizes, silent. An earlier draft of mine claimed ``20 minutes against 20
hours''; the 20 minutes described one cell and the 20 hours had no source. The figure shows the
measured totals.

\subsection{Multiple comparisons}\label{sec:multiplicity}
The first battery ran 19 formal tests: nine divergence comparisons (three build pairs in three
domains), six HellaSwag pairs, two HumanEval+ pairs, the RULER closure test and the draft-depth sign
test. All were two-sided except the sign test, and none was pre-registered except the RULER analysis,
so the family is exploratory. Under a Holm correction, which tolerates dependence between tests, the
nine divergence comparisons pass whether computed from the tool's token-level errors or from the paired
window-level analysis of \cref{tab:paired-kld} (smallest paired $t$ is 4.1 on 31 degrees of freedom,
$p\approx3\times10^{-4}$, against a Holm threshold near 0.005), and so does the RULER closure
difference. The draft-depth sign test ($p=0.0059$) fails Holm at its rank; its eleven comparisons also
share cells, so they are not independent, and I treat it as directional only. No task-benchmark
comparison passes. Given their observed discordance, six of those eight could not have returned
$p<0.05$.
